% pracovni_list.tex -- Lekce 05: Náhoda a hry
\newcommand{\lekce}{Lekce 05 -- Pracovní list}
% preamble-pl.tex -- sdílená preamble pro pracovní listy
% Includuje se z ucitel/lekceXX/pracovni_list.tex jako:
%   % preamble-pl.tex -- sdílená preamble pro pracovní listy
% Includuje se z ucitel/lekceXX/pracovni_list.tex jako:
%   % preamble-pl.tex -- sdílená preamble pro pracovní listy
% Includuje se z ucitel/lekceXX/pracovni_list.tex jako:
%   \input{../spolecne/preamble-pl.tex}

\documentclass[a4paper,12pt]{article}

% --- Jazyk a font ---
\usepackage{polyglossia}
\setmainlanguage{czech}
\usepackage{fontspec}

% --- Okraje ---
\usepackage[top=20mm, bottom=20mm, left=22mm, right=22mm]{geometry}

% --- Česky korektní uvozovky ---
\usepackage{csquotes}
\newcommand{\uv}[1]{\enquote{#1}}

% --- Typografie ---
\usepackage{microtype}
\usepackage{parskip}          % odstavce odděleny mezerou, ne odsazením

% --- Barvy a grafika ---
\usepackage[table]{xcolor}
\definecolor{microbit-blue}{HTML}{1E4D8C}
\definecolor{code-bg}{HTML}{F5F5F5}
\definecolor{task-bg}{HTML}{EEF4FF}

% --- Tabulky ---
\usepackage{booktabs}
\usepackage{multirow}

% --- Zdrojový kód (Python) ---
\usepackage{listings}
\lstset{
  language=Python,
  basicstyle=\ttfamily\small,
  backgroundcolor=\color{code-bg},
  frame=single,
  framerule=0pt,
  rulecolor=\color{gray!30},
  xleftmargin=6pt,
  xrightmargin=6pt,
  breaklines=true,
  showstringspaces=false,
  keywordstyle=\color{microbit-blue}\bfseries,
  stringstyle=\color{teal},
  commentstyle=\color{gray}\itshape,
  literate=
    {á}{{\'a}}1 {é}{{\'e}}1 {í}{{\'i}}1 {ó}{{\'o}}1 {ú}{{\'u}}1
    {ů}{{\r{u}}}1 {č}{{\v{c}}}1 {š}{{\v{s}}}1 {ž}{{\v{z}}}1
    {Á}{{\'A}}1 {É}{{\'E}}1 {Í}{{\'I}}1 {Ó}{{\'O}}1 {Ú}{{\'U}}1
    {Č}{{\v{C}}}1 {Š}{{\v{S}}}1 {Ž}{{\v{Z}}}1 {ň}{{\v{n}}}1 {ř}{{\v{r}}}1,
}

% --- Zadání úkolů ---
\usepackage[most]{tcolorbox}
\tcbuselibrary{breakable}

\newcounter{ukol}[section]
\newenvironment{ukol}[1][]{%
  \stepcounter{ukol}%
  \begin{tcolorbox}[
    colback=task-bg,
    colframe=microbit-blue,
    fonttitle=\bfseries,
    title={Úkol~\theukol\ifx\relax#1\relax\else{: #1}\fi},
    breakable,
    left=6pt, right=6pt, top=4pt, bottom=4pt
  ]%
}{%
  \end{tcolorbox}%
}

% Inline kód — underscore package zajistí, ze _ funguje v texttt
\usepackage{underscore}
\newcommand{\pyinline}[1]{\texttt{#1}}
% Pojem (zvýraznění termínu) — stejné jako v preamble-prez.tex
\newcommand{\pojem}[1]{\textbf{#1}}

% Místo pro odpověď studenta
\newcommand{\odpoved}[1][3cm]{%
  \vspace{#1}%
}

% --- Záhlaví / zápatí ---
\usepackage{fancyhdr}
\pagestyle{fancy}
\fancyhf{}
\renewcommand{\headrulewidth}{0.4pt}
\fancyhead[L]{\small\textcolor{microbit-blue}{\textbf{Úvod do Pythonu na Micro:bitu}}}
\fancyhead[R]{\small\textcolor{gray}{\lekce}}
\fancyfoot[C]{\small\thepage}

% --- Ostatní ---
\usepackage{enumitem}
\usepackage{graphicx}
\usepackage{hyperref}
\hypersetup{
  colorlinks=true,
  linkcolor=microbit-blue,
  urlcolor=microbit-blue,
}


\documentclass[a4paper,12pt]{article}

% --- Jazyk a font ---
\usepackage{polyglossia}
\setmainlanguage{czech}
\usepackage{fontspec}

% --- Okraje ---
\usepackage[top=20mm, bottom=20mm, left=22mm, right=22mm]{geometry}

% --- Česky korektní uvozovky ---
\usepackage{csquotes}
\newcommand{\uv}[1]{\enquote{#1}}

% --- Typografie ---
\usepackage{microtype}
\usepackage{parskip}          % odstavce odděleny mezerou, ne odsazením

% --- Barvy a grafika ---
\usepackage[table]{xcolor}
\definecolor{microbit-blue}{HTML}{1E4D8C}
\definecolor{code-bg}{HTML}{F5F5F5}
\definecolor{task-bg}{HTML}{EEF4FF}

% --- Tabulky ---
\usepackage{booktabs}
\usepackage{multirow}

% --- Zdrojový kód (Python) ---
\usepackage{listings}
\lstset{
  language=Python,
  basicstyle=\ttfamily\small,
  backgroundcolor=\color{code-bg},
  frame=single,
  framerule=0pt,
  rulecolor=\color{gray!30},
  xleftmargin=6pt,
  xrightmargin=6pt,
  breaklines=true,
  showstringspaces=false,
  keywordstyle=\color{microbit-blue}\bfseries,
  stringstyle=\color{teal},
  commentstyle=\color{gray}\itshape,
  literate=
    {á}{{\'a}}1 {é}{{\'e}}1 {í}{{\'i}}1 {ó}{{\'o}}1 {ú}{{\'u}}1
    {ů}{{\r{u}}}1 {č}{{\v{c}}}1 {š}{{\v{s}}}1 {ž}{{\v{z}}}1
    {Á}{{\'A}}1 {É}{{\'E}}1 {Í}{{\'I}}1 {Ó}{{\'O}}1 {Ú}{{\'U}}1
    {Č}{{\v{C}}}1 {Š}{{\v{S}}}1 {Ž}{{\v{Z}}}1 {ň}{{\v{n}}}1 {ř}{{\v{r}}}1,
}

% --- Zadání úkolů ---
\usepackage[most]{tcolorbox}
\tcbuselibrary{breakable}

\newcounter{ukol}[section]
\newenvironment{ukol}[1][]{%
  \stepcounter{ukol}%
  \begin{tcolorbox}[
    colback=task-bg,
    colframe=microbit-blue,
    fonttitle=\bfseries,
    title={Úkol~\theukol\ifx\relax#1\relax\else{: #1}\fi},
    breakable,
    left=6pt, right=6pt, top=4pt, bottom=4pt
  ]%
}{%
  \end{tcolorbox}%
}

% Inline kód — underscore package zajistí, ze _ funguje v texttt
\usepackage{underscore}
\newcommand{\pyinline}[1]{\texttt{#1}}
% Pojem (zvýraznění termínu) — stejné jako v preamble-prez.tex
\newcommand{\pojem}[1]{\textbf{#1}}

% Místo pro odpověď studenta
\newcommand{\odpoved}[1][3cm]{%
  \vspace{#1}%
}

% --- Záhlaví / zápatí ---
\usepackage{fancyhdr}
\pagestyle{fancy}
\fancyhf{}
\renewcommand{\headrulewidth}{0.4pt}
\fancyhead[L]{\small\textcolor{microbit-blue}{\textbf{Úvod do Pythonu na Micro:bitu}}}
\fancyhead[R]{\small\textcolor{gray}{\lekce}}
\fancyfoot[C]{\small\thepage}

% --- Ostatní ---
\usepackage{enumitem}
\usepackage{graphicx}
\usepackage{hyperref}
\hypersetup{
  colorlinks=true,
  linkcolor=microbit-blue,
  urlcolor=microbit-blue,
}


\documentclass[a4paper,12pt]{article}

% --- Jazyk a font ---
\usepackage{polyglossia}
\setmainlanguage{czech}
\usepackage{fontspec}

% --- Okraje ---
\usepackage[top=20mm, bottom=20mm, left=22mm, right=22mm]{geometry}

% --- Česky korektní uvozovky ---
\usepackage{csquotes}
\newcommand{\uv}[1]{\enquote{#1}}

% --- Typografie ---
\usepackage{microtype}
\usepackage{parskip}          % odstavce odděleny mezerou, ne odsazením

% --- Barvy a grafika ---
\usepackage[table]{xcolor}
\definecolor{microbit-blue}{HTML}{1E4D8C}
\definecolor{code-bg}{HTML}{F5F5F5}
\definecolor{task-bg}{HTML}{EEF4FF}

% --- Tabulky ---
\usepackage{booktabs}
\usepackage{multirow}

% --- Zdrojový kód (Python) ---
\usepackage{listings}
\lstset{
  language=Python,
  basicstyle=\ttfamily\small,
  backgroundcolor=\color{code-bg},
  frame=single,
  framerule=0pt,
  rulecolor=\color{gray!30},
  xleftmargin=6pt,
  xrightmargin=6pt,
  breaklines=true,
  showstringspaces=false,
  keywordstyle=\color{microbit-blue}\bfseries,
  stringstyle=\color{teal},
  commentstyle=\color{gray}\itshape,
  literate=
    {á}{{\'a}}1 {é}{{\'e}}1 {í}{{\'i}}1 {ó}{{\'o}}1 {ú}{{\'u}}1
    {ů}{{\r{u}}}1 {č}{{\v{c}}}1 {š}{{\v{s}}}1 {ž}{{\v{z}}}1
    {Á}{{\'A}}1 {É}{{\'E}}1 {Í}{{\'I}}1 {Ó}{{\'O}}1 {Ú}{{\'U}}1
    {Č}{{\v{C}}}1 {Š}{{\v{S}}}1 {Ž}{{\v{Z}}}1 {ň}{{\v{n}}}1 {ř}{{\v{r}}}1,
}

% --- Zadání úkolů ---
\usepackage[most]{tcolorbox}
\tcbuselibrary{breakable}

\newcounter{ukol}[section]
\newenvironment{ukol}[1][]{%
  \stepcounter{ukol}%
  \begin{tcolorbox}[
    colback=task-bg,
    colframe=microbit-blue,
    fonttitle=\bfseries,
    title={Úkol~\theukol\ifx\relax#1\relax\else{: #1}\fi},
    breakable,
    left=6pt, right=6pt, top=4pt, bottom=4pt
  ]%
}{%
  \end{tcolorbox}%
}

% Inline kód — underscore package zajistí, ze _ funguje v texttt
\usepackage{underscore}
\newcommand{\pyinline}[1]{\texttt{#1}}
% Pojem (zvýraznění termínu) — stejné jako v preamble-prez.tex
\newcommand{\pojem}[1]{\textbf{#1}}

% Místo pro odpověď studenta
\newcommand{\odpoved}[1][3cm]{%
  \vspace{#1}%
}

% --- Záhlaví / zápatí ---
\usepackage{fancyhdr}
\pagestyle{fancy}
\fancyhf{}
\renewcommand{\headrulewidth}{0.4pt}
\fancyhead[L]{\small\textcolor{microbit-blue}{\textbf{Úvod do Pythonu na Micro:bitu}}}
\fancyhead[R]{\small\textcolor{gray}{\lekce}}
\fancyfoot[C]{\small\thepage}

% --- Ostatní ---
\usepackage{enumitem}
\usepackage{graphicx}
\usepackage{hyperref}
\hypersetup{
  colorlinks=true,
  linkcolor=microbit-blue,
  urlcolor=microbit-blue,
}


\begin{document}

\begin{center}
  {\Large\bfseries\textcolor{microbit-blue}{Lekce 05: Náhoda a hry}}\\[4pt]
  {\large Pracovní list}
\end{center}

\medskip\noindent
Python umí generovat náhodná čísla pomocí modulu \pyinline{random}.
Dnes ho použijeme spolu s gestem třepání micro:bitu na vytvoření her.

% ------------------------------------------------------------------
\begin{ukol}[Hod kostkou]

Zatřes micro:bitem — hodíš kostkou.

\begin{lstlisting}
from microbit import *
import random

while True:
    if accelerometer.was_gesture("shake"):
        cislo = random.randint(1, 6)
        display.show(str(cislo))
    sleep(100)
\end{lstlisting}

\hadej{Hod kostkou — předem tipni 5 výsledků, pak porovnej:

\medskip
\begin{tabular}{llllll}
  & Hod 1 & Hod 2 & Hod 3 & Hod 4 & Hod 5 \\
  \midrule
  Odhad & \hspace{1cm} & \hspace{1cm} & \hspace{1cm} & \hspace{1cm} & \hspace{1cm} \\[4pt]
  Skutečnost & & & & & \\
\end{tabular}}

\begin{enumerate}
  \item Zadej program a spusť ho. Zatřes micro:bitem — zobrazí se číslo 1–6.
  \item Co dělá \pyinline{random.randint(1, 6)}?
        Může vrátit číslo 6? A číslo 0?
  \item Uprav program na desetistěnnou kostku (1–10).
  \item Co dělá \pyinline{sleep(100)} na konci smyčky?
        Zkus ho odstranit — co se změní a proč?
\end{enumerate}

\odpoved[2.5cm]

\begin{tcolorbox}[colback=code-bg, colframe=gray!40,
                  fonttitle=\small\bfseries, title=import a randint]
  \small
  \begin{itemize}[leftmargin=*]
    \item \pyinline{import random} — načte modul s funkcemi pro náhodu.
    \item \pyinline{random.randint(a, b)} — vrátí náhodné celé číslo
          od \pyinline{a} do \pyinline{b} \textbf{včetně}.
    \item \pyinline{accelerometer.was_gesture("shake")} — vrátí \pyinline{True},
          pokud micro:bit zaznamenal třepání od posledního volání.
  \end{itemize}
\end{tcolorbox}

\end{ukol}

% ------------------------------------------------------------------
\begin{ukol}[Náhodný obrázek]

Místo čísla zobrazíme náhodný obrázek ze seznamu:

\begin{lstlisting}
from microbit import *
import random

obrazky = [Image.HAPPY, Image.SAD, Image.SURPRISED,
           Image.ANGRY, Image.CONFUSED, Image.ASLEEP]

while True:
    if accelerometer.was_gesture("shake"):
        display.show(random.choice(obrazky))
    sleep(100)
\end{lstlisting}

\begin{enumerate}
  \item Co dělá \pyinline{random.choice(obrazky)}?
        Jak se liší od \pyinline{random.randint()}?
  \item Přidej do seznamu \pyinline{obrazky} další obrázky dle vlastního výběru.
  \item Uprav program tak, aby se spolu s obrázkem také skroloval krátký text
        (např. „Dnes bude hezky!").
\end{enumerate}

\odpoved[2cm]

\end{ukol}

% ------------------------------------------------------------------
\begin{ukol}[Věštírna]

Vytvoř program, který po zatřesení zobrazí náhodnou odpověď z vlastního seznamu.
Věštírna může předpovídat cokoliv — odpovědi Ano/Ne, číslo štěstí, předpověď
počasí, náladu\ldots

\medskip
Podmínky:
\begin{itemize}
  \item seznam musí mít alespoň 4 různé možnosti
  \item použij \pyinline{random.choice()}
  \item pro delší texty použij \pyinline{display.scroll()},
        pro obrázky nebo jednoznakové výstupy \pyinline{display.show()}
\end{itemize}

\odpoved[4.5cm]

\noindent Výsledný program ulož a připrav k odevzdání do Google Classroom.

\end{ukol}

% ------------------------------------------------------------------
\begin{bonusukol}[Uhádni číslo — hra pro jednoho]

Micro:bit si vybere tajné číslo 1–3. Ty hádáš: tlačítko A = hádám 1,
tlačítko B = hádám 2, obě = hádám 3.

\begin{lstlisting}
from microbit import *
import random

tajne = random.randint(1, 3)

while True:
    # is_pressed() pro A+B — was_pressed() by se po prvním volání resetovalo
    if button_a.is_pressed() and button_b.is_pressed():
        tip = 3
    elif button_a.was_pressed():
        tip = 1
    elif button_b.was_pressed():
        tip = 2
    else:
        tip = 0

    if tip != 0:
        if tip == tajne:
            display.show(Image.HAPPY)
        else:
            display.show(Image.SAD)
            display.scroll(str(tajne))   # prozrad spravnou odpoved
        tajne = random.randint(1, 3)     # nova hra
    sleep(50)
\end{lstlisting}

\begin{enumerate}
  \item Zadej program a zahraj si — uhodneš správné číslo?
  \item Proč se \pyinline{tajne = random.randint(1, 3)} opakuje i uvnitř smyčky?
        Co by se stalo bez tohoto řádku?
  \item \textbf{Výzva:} Přidej proměnnou \pyinline{skore}, která počítá správné tipy.
        Po každé správné odpovědi zobraz aktuální skóre.
\end{enumerate}

\odpoved[2.5cm]

\end{bonusukol}

\end{document}
